\documentclass[12pt]{article}
\usepackage{amsmath}
\usepackage{amssymb}  % Provides \mathbb and other symbols
\usepackage{geometry}
\usepackage{graphicx}
\usepackage{subcaption}
\usepackage{enumitem}
\usepackage{listings}
\usepackage{hyperref}
\usepackage{ulem}
\usepackage{xcolor}
\usepackage{amsthm}
\usepackage{tcolorbox}  % For colored boxes
\usepackage[table]{xcolor} % For coloring table cells
\usepackage{array}         % For better table alignment
\usepackage{mdframed}
\usepackage{cancel}
\usepackage{amsmath,amssymb}

% Set the top margin
\geometry{top=1in} % You can adjust the measurement as needed


\begin{document}

\begin{titlepage}
\centering
\vspace*{\fill} % Adds vertical space to center the title block vertically

{\Huge Theoretical Physics for Mathematicians} \\[0.5cm] % Title 2
{\Huge Home Work 2}\\[1.5cm] % Title 3
{\Large Student: 26-712-224}\\[0.5cm] % Student ID
{\Large \today} % Date, \today adds today's date

\vspace*{\fill}
\end{titlepage}
\newpage
\noindent

\section*{Question 1 - \(\star\) Hamilton's Principle: The Euler--Lagrange equations and their invariances}
Let \(L=L(q,\dot q;t)\) be smooth and let
\begin{align*}
S[q]=\int_{t_0}^{t_1}L(q(t),\dot q(t);t)\,dt.
\end{align*}

\begin{enumerate}
\item[(a)] \textbf{(Fundamental lemma).}
Let \(F:[t_0,t_1]\to\mathbb{R}^f\) be continuous and suppose
\begin{align*}
\int_{t_0}^{t_1}F_a(t)\eta^a(t)\,dt=0
\end{align*}
for every smooth \(\eta\) with compact support in \((t_0,t_1)\).
Prove that \(F=0\).
\textit{Hint: if \(F_a(t_*)\neq 0\) for some \(t_*\) and some index \(a\),
construct a bump function concentrated near \(t_*\).}

\item[(b)] Compute the first variation
\begin{align*}
\delta S[q](\eta)
=
\left.
\frac{d}{d\epsilon}S[q+\epsilon\eta]
\right|_{\epsilon=0}
\end{align*}
for variations with \(\eta(t_0)=\eta(t_1)=0\), and deduce from (a) that
\(\delta S=0\) is equivalent to the Euler--Lagrange equations.

\item[(c)] Let \(F=F(q,t)\) be smooth and set
\begin{align*}
L'
=
L+\frac{d}{dt}F(q,t),
\qquad\text{i.e.}\qquad
L'
=
L+\frac{\partial F}{\partial q^a}\dot q^a
+\frac{\partial F}{\partial t}.
\end{align*}
Show that \(L\) and \(L'\) have exactly the same Euler--Lagrange equations.
How do the actions \(S\) and \(S'\) differ?

\item[(d)] \textbf{(Coordinate invariance).}
Let \(q=q(Q)\) be a diffeomorphic change of generalized coordinates and let
\begin{align*}
\widetilde L(Q,\dot Q;t)
=
L\left(q(Q),\frac{\partial q}{\partial Q}\dot Q;t\right).
\end{align*}
Prove the identity
\begin{align*}
\frac{d}{dt}\frac{\partial\widetilde L}{\partial \dot Q^b}
-
\frac{\partial\widetilde L}{\partial Q^b}
=
\frac{\partial q^a}{\partial Q^b}
\left(
\frac{d}{dt}\frac{\partial L}{\partial\dot q^a}
-
\frac{\partial L}{\partial q^a}
\right),
\end{align*}
and conclude that the Euler--Lagrange equations are coordinate independent.
Why does the Cartesian component form \(m\ddot x^a=F^a\) acquire connection
terms in curvilinear coordinates?
\end{enumerate}

\newpage
\noindent
Answer: \\
\\

\newpage
\noindent
\section*{Question 1 - \(\star\) Oscillator and pendulum, with and without multipliers}
\begin{enumerate}
\item[(a)] For
\begin{align*}
L=\frac12 m\dot q^2-\frac12 kq^2
\end{align*}
derive the equation of motion and its general solution, compute
\begin{align*}
E_L
=
\dot q\frac{\partial L}{\partial\dot q}-L
\end{align*}
and verify directly that \(\dot E_L=0\) on solutions of the Euler--Lagrange
equation.

\item[(b)] For the plane pendulum of length \(\ell\) use the angle \(\theta\)
(measured from the downward vertical) to derive
\begin{align*}
\ddot\theta+\frac{g}{\ell}\sin\theta=0.
\end{align*}
Determine the small-oscillation frequency, and show that energy conservation
reduces the problem to the quadrature
\begin{align*}
\dot\theta^2
=
\frac{2g}{\ell}
\left(
\cos\theta-\cos\theta_0
\right),
\end{align*}
where \(0<\theta_0<\pi\) is the amplitude of a libration.
Deduce the exact period as an integral and check that it tends to
\(2\pi\sqrt{\ell/g}\) as \(\theta_0\to0\).

\item[(c)] Treat the same pendulum in \(\mathbb{R}^2\) with the constraint
\begin{align*}
G(x,y)=x^2+y^2-\ell^2=0
\end{align*}
and
\begin{align*}
L=\frac12 m(\dot x^2+\dot y^2)-mgy.
\end{align*}
Write down the multiplier equations, and determine \(\lambda\) explicitly
as a function of \((x,y,\dot x,\dot y)\).

\item[(d)] Write the constraint force as
\begin{align*}
\vec F_{\mathrm{constraint}}
=
-\mathcal{T}\,\hat e_r,
\end{align*}
where
\begin{align*}
\hat e_r=(\sin\theta,-\cos\theta)
\end{align*}
points outwards.
Show that the signed inward tension is
\begin{align*}
\mathcal{T}
=
m\left(g\cos\theta+\ell\dot\theta^2\right),
\qquad
\left|\vec F_{\mathrm{constraint}}\right|
=
|\mathcal{T}|.
\end{align*}

Explain the distinction between a rigid rod, which can exert compression,
and a taut string.
Confirm that (c) reproduces the equation of motion of (b).
\end{enumerate}

\newpage
\noindent
Answer: \\
\\

\end{document}